\documentclass[12pt]{article}

\usepackage{amsmath,amssymb,graphicx,hyperref,tikz}
\usepackage{geometry}
\geometry{margin=1in}

\title{Geometric-Topological Derivation of the Fine-Structure Constant via TTT 6D Polyhedral Dynamics}
\author{Shinketsu}
\date{\today}

\begin{document}
\maketitle

\begin{abstract}
The fine-structure constant $\alpha^{-1} \approx 137.035999$ remains theoretically unexplained in the Standard Model.
We propose a geometric-topological derivation based on the TTT 6D polyhedral transition model.
The constant emerges as the fixed point of an infinite renormalization process.
\end{abstract}

\section{Introduction}
% --- あなたの本文（完成版）をここに貼る ---

\section{TTT 6D State Model}
% --- 本文 ---

\section{Polyhedral Transition Chain}

\begin{figure}[htbp]
\centering
\includegraphics[width=0.9\textwidth]{poly_transition.png}
\caption{TTT polyhedral transition chain forming a closed orbit with winding number $W=3$.}
\label{fig:polyhedral_transition}
\end{figure}

\section{Convergence of the Fine-Structure Constant}

\begin{figure}[htbp]
\centering
\includegraphics[width=0.75\textwidth]{alpha_convergence.png}
\caption{Convergence of $\alpha^{-1}$ in the TTT model. The value approaches the experimental constant $137.035999$.}
\label{fig:alpha_convergence}
\end{figure}

\section{Conclusion}
% --- 本文 ---

\appendix
\section{Appendix: Python Code}

\begin{verbatim}
% Appendix の Python コード全文をここに貼る
\end{verbatim}

\section*{References}
\begin{thebibliography}{99}
% --- 完成版 References をここに貼る ---
\end{thebibliography}

\end{document}
